\chapter*{}
\noindent
\phantomsection{\MakeUppercase{\textbf{Conclusões}}}
\addcontentsline{toc}{chapter}{CONCLUSÕES}
\newline
\newline

O objetivo geral deste trabalho foi alcançado. Foi desenvolvido um
sistema de monitoramento e alerta para barragens de rejeitos de
mineração, baseado em IoT, que combina a saturação do solo, obtida no
local, com a chuva observada e prevista, e emite alertas em três
níveis, de forma local e remota, com alimentação redundante por bateria.

Os seis objetivos específicos do \autoref{cap:objetivos} foram
atendidos (\autoref{tab:cobertura_objetivos}). Nos 31 casos de teste
executados, a saída observada foi igual à saída esperada
(\autoref{tab:res-sintese}).

O modelo de alerta aplica um limiar hidrometeorológico: o nível sobe
quando a causa, o solo úmido, e o gatilho, a chuva, ocorrem em conjunto
\cite{bogaard2018hydrological}. A saturação relativa $S$ é obtida pelo
higrômetro, e a chuva $P_{72}$ soma 48~horas observadas na BNDMET e
24~horas previstas pelo OpenWeatherMap \cite[p.~5]{mirus2018developing}.

Cada limiar do modelo tem origem na literatura
(\autoref{tab:limiares}). Os de saturação, $S_1 = 0{,}64$ e
$S_2 = 0{,}86$, vêm de \citeonline[p.~9, Tabela~1]{mirus2018developing},
e os de chuva, $P_1 = 60$\,mm e $P_2 = 100$\,mm em 72~horas, de
\citeonline[p.~74, Tabela~3]{mendes2020proposicao}.

O alerta local é feito pelos LEDs de cada nível e pelo
\textit{buzzer}, acionado no VERMELHO. Como o nível é calculado no
próprio dispositivo, esse alerta não depende da rede nem do
\textit{backend}, e as regras de contingência o mantêm quando falta o
higrômetro ou os dados de chuva (\autoref{tab:contingencia}).

O alerta remoto é feito pela plataforma \textit{web}. O morador se
cadastra na página inicial, e o administrador acompanha o painel e
envia o e-mail de alerta verde, amarelo ou vermelho, com a orientação
sobre a necessidade de deslocamento (\autoref{qua:mensagens-alerta}).
No teste, o e-mail foi enviado aos cinco moradores
cadastrados (\autoref{fig:res-w17b}), e o recebimento foi conferido na
caixa de entrada de um deles (\autoref{fig:res-w19}).

A redundância de energia foi dimensionada por cálculo. Para 24~horas
de supervisão e 15~minutos de alarme \cite[item~5.3]{cbmmg2017it14},
com margem de 20\% \cite{hammerberg2021calculation}, a bateria precisa
de 7816\,mAh (\autoref{eq:cap-celula}). Por isso, foram adotadas três
células NCR18650B em paralelo, com 9600\,mAh \cite{panasonic_ncr18650b}.

Na simulação, com a falta de rede, o dispositivo passou à bateria e
continuou a calcular o nível e a enviar as leituras
(\autoref{fig:res-e02}). Os avisos de bateria baixa, bateria crítica e
falha de recarga também foram demonstrados (\autoref{tab:res-energia}).
A autonomia teórica em supervisão é de cerca de 36~horas
(\autoref{tab:autonomia}).

A Política Nacional de Segurança de Barragens prevê, no plano de
emergência, rotas de fuga e pontos de encontro sinalizados e um sistema
sonoro ou outra solução tecnológica para situações de alerta
\cite{brasil2010pnsb,brasil2020lei14066}. O sistema desenvolvido segue
essa lógica, com o alerta sonoro local e o e-mail que orienta o
deslocamento ao ponto de encontro.

O sistema, porém, é complementar: ele não substitui a sirene, o plano de
emergência da barragem nem a Defesa Civil, e as suas cores não
correspondem aos níveis de emergência definidos pela
\citeonline{anm2022resolucao95}.

Os resultados se limitam, como discutido na \autoref{sec:limitacoes},
a um único ponto de medição, com um sensor que não mede a umidade
volumétrica, limiares de referência não calibrados para Brumadinho nem
para rejeito de minério de ferro e uma bateria apenas calculada e
simulada, em testes acelerados.

A partir dessas limitações, são propostos os seguintes trabalhos
futuros:

\begin{enumerate}

  \item calibrar os limiares $S_1$, $S_2$, $P_1$ e $P_2$ para o local
  monitorado e para o rejeito de minério de ferro;

  \item substituir o higrômetro FC-28 por um sensor de umidade
  volumétrica, para obter $S$ na mesma escala em que $S_1$ e $S_2$
  foram definidos \cite[p.~5]{mirus2018developing};

  \item montar a bateria, o carregador e a detecção da rede no
  protótipo físico e ensaiar a autonomia real frente ao requisito de
  24~horas de supervisão e 15~minutos de alarme
  \cite[item~5.3]{cbmmg2017it14};

  \item instalar múltiplos pontos de medição, em posições e
  profundidades diferentes da estrutura;

  \item complementar os dias sem registro na estação da BNDMET com os
  dados de uma segunda estação próxima.

\end{enumerate}